% Generated draft: requires mathematical and source review.
\documentclass[11pt]{article}
\usepackage[margin=25mm]{geometry}
\usepackage{fontspec}
\setmainfont{Latin Modern Roman}
\usepackage{amsmath,amssymb,mathtools,unicode-math}
\setmathfont{Latin Modern Math}
\usepackage{enumitem,needspace,xurl,hyperref}
\setlength{\parindent}{0pt}
\setlength{\parskip}{5pt}
\setlength{\emergencystretch}{3em}
\newcommand{\N}{\mathbb{N}}
\newcommand{\Z}{\mathbb{Z}}
\newcommand{\Q}{\mathbb{Q}}
\newcommand{\R}{\mathbb{R}}
\newcommand{\C}{\mathbb{C}}
\newcommand{\F}{\mathbb{F}}
\newcommand{\A}{\mathbb{A}}
\newcommand{\PP}{\mathbb{P}}
\newcommand{\E}{\mathbb{E}}
\newcommand{\eps}{\varepsilon}
\newcommand{\dd}{\,\mathrm{d}}
\begin{document}
\section*{UnsolvedMath ID: 6200092}
\textbf{Source problem code:} catalogue entry
\subsection{Definitions and mathematical statement}
A CAT(0) space is a geodesic metric space whose geodesic triangles satisfy the CAT(0) comparison inequality with Euclidean triangles. Consider an isometric action of a group \(G\) on such a space \(X\).

For comparison, a subset \(Q\) of a geodesic metric space is quasiconvex if some \(R\ge0\) satisfies
\[
 [x,y]\subseteq\{z:d(z,Q)\le R\}\qquad(x,y\in Q),
\]
for every joining geodesic. Applied to an orbit this gives one familiar action-level condition. Another familiar condition is a proper action cocompact on a suitable invariant convex subset, often called convex cocompactness. These are background notions, not a definition of geometric finiteness in all CAT(0) spaces.

Kapovich's problem is to develop a theory of geometric finiteness for group actions on CAT(0) spaces, extending the role of quasiconvex actions on Gromov-hyperbolic spaces and convex-cocompact real-hyperbolic isometry groups. The source supplies no proposed definition, characterization theorem, or exact axioms for the desired extension. It does not fix properness or dimension of \(X\), action properness, finite generation, allowed parabolic behavior, or boundary conditions. Consequently the target is a research program, not a fully quantified assertion that one of the displayed familiar conditions is the universal answer.

The supplied summary reports progress for settings such as rank-one actions and cube complexes but no complete general theory. It provides no precise examples or theorem citations. These qualified partial-progress reports and the 2005 problem label are retained without an independent current-status assessment.
\subsection{Short English statement}
Develop a useful notion and theory of geometric finiteness for CAT(0) group actions, extending familiar hyperbolic-space theories. The source does not specify the intended axioms, action class, or a single theorem to prove.
\subsection{Sources}
\begin{itemize}
\item Local registry: \url{data/unsolvedmath/problems.json}, permanent UnsolvedMath ID \(6200092\), source code catalogue entry; supplied statement, background, and research summary, updated 2026-08-06.
\item Misha Kapovich, \emph{Problems on Boundaries of Groups and Kleinian Groups} (2005), Problem 92, PDF page 23: \url{https://www.math.ucdavis.edu/~kapovich/EPR/problems.pdf}. Supplied normalized PDF-text extraction accessed 2026-07-29; current status marked as needing review.
\item The summary mentions rank-one and cube-complex progress without identifiable original-paper references or exact statements.
\end{itemize}
\end{document}
